\section{Einleitung}
\label{sec:einleitung}

Soziale Medien haben sich in den vergangenen Jahren zu einer bedeutenden Informationsquelle für Finanzmärkte entwickelt. Plattformen wie X\,(ehem.\ Twitter) ermöglichen es, öffentlich zugängliche Stimmungsäußerungen von Marktteilnehmern in nahezu Echtzeit zu beobachten. Gleichzeitig sind Aktienkurse von Informationen und kollektiven Erwartungen abhängig, die sich über digitale Kanäle rasch verbreiten. Diese Arbeit untersucht daher die Frage, ob und in welchem Ausmaß das auf X\,(Twitter) gemessene Sentiment von Marktteilnehmern mit den Kursrenditen ausgewählter Unternehmen zusammenhängt.

\subsection*{Forschungsfrage}

Die zentrale Forschungsfrage lautet:

\begin{quote}
\textit{Lässt sich aus dem täglichen Sentiment auf X\,(Twitter) ein statistisch signifikanter Zusammenhang mit den Tagesrenditen börsennotierter Unternehmen aus den Bereichen Energie, Elektromobilität und Wärmepumpen ableiten?}
\end{quote}

Konkret wird untersucht, ob Sentiment-Signale den Kursrenditen vorausgehen, ihnen folgen oder ob kein erklärender Zusammenhang besteht. Dabei wird auch geprüft, ob vergangenes Sentiment im Sinne des Granger-Kausalitätstests Vorhersageinformation über zukünftige Renditen enthält.

\subsection*{Fokus und Relevanz}

Das Untersuchungsfeld umfasst zwölf Unternehmen aus drei thematischen Segmenten: etablierte Energiekonzerne (BP, ExxonMobil, Shell, TotalEnergies, Eni), Hersteller von Wärmepumpen und Klimatisierungstechnik (Daikin, NIBE, Carrier Global, Mitsubishi Electric) sowie Automobilhersteller mit Elektromobilitätsfokus (Tesla, BYD, Volkswagen, Hyundai). Diese Auswahl ist relevant, weil sie Unternehmen umfasst, die von laufenden Diskursen über Energiewende und Klimapolitik betroffen sind und daher eine erhöhte Exposition gegenüber öffentlichem Sentiment erwarten lassen.

\subsection*{Beitrag und Aufbau}

Der Beitrag dieser Arbeit liegt in der Verknüpfung eines aktuellen, selbst erhobenen Social-Media-Datensatzes mit historischen Börsenkursdaten sowie in der systematischen Anwendung mehrerer quantitativer Methoden auf diesen spezifischen Ausschnitt. Abschnitt~\ref{sec:methodik} beschreibt die Datenerhebung, Datenaufbereitung und die eingesetzten Analysemethoden. Abschnitt~\ref{sec:ergebnisse} stellt die Ergebnisse dar. Abschnitt~\ref{sec:diskussion} bewertet diese kritisch, geht auf Limitationen ein und ordnet die Befunde in den Forschungskontext ein. Abschnitt~\ref{sec:fazit} schließt mit einem Fazit.
